\chapter{Rash}

\section*{Definition}
Abnormal change in skin color or texture possibly involving bumps, patches, blisters.

\section*{What Happens in Your Body}
Numerous etiations including allergic reactions, infections, autoimmune diseases, drug eruptions, mechanical factors.

\section*{Causes}
\begin{itemize}
  \item Allergic contact dermatitis
  \item Viral exanthems
  \item Drug eruption
  \item Autoimmune diseases
  \item Fungal infections
\end{itemize}

\section*{When to See a Doctor}
See your doctor if:
\begin{itemize}
  \item Symptoms are severe or getting worse over time
  \item Evaluate with systemic features such as fever joint — seek medical evaluation
  \item You have other concerning symptoms such as fever, weight loss, or bleeding
\end{itemize}

\section*{Related Symptoms}
\begin{itemize}
  \item Itching
  \item Pain
  \item Blisters/scaling
  \item Redness
  \item Swelling
\end{itemize}
\textbf{Important:} If this symptom persists, your doctor will check for serious underlying causes.

\section*{Self-Care / Home Management}
\begin{itemize}
  \item Keep the affected area clean and dry; wash gently with mild, fragrance-free soap and lukewarm water.
  \item Avoid scratching — use cool compresses, calamine lotion, or OTC hydrocortisone cream to relieve itching.
  \item Apply moisturizer regularly to maintain skin barrier integrity, especially after bathing.
  \item Identify and avoid potential triggers (irritants, allergens, excessive heat, tight clothing).
  \item Seek medical evaluation if the rash spreads rapidly, develops blisters, shows signs of infection (pus, increasing redness, warmth), or is accompanied by fever.
\end{itemize}

\section*{How Your Doctor Will Evaluate You}
\begin{itemize}
  \item History includes onset, location, progression, associated symptoms (itching, pain, fever), exposure to triggers, and prior skin conditions.
  \item Physical examination assesses morphology (macules, papules, vesicles, etc.), distribution, and extent of skin involvement.
  \item Diagnostic tests may include skin scraping for fungal examination, patch testing for allergic contact dermatitis, or biopsy for suspicious lesions.
  \item Laboratory studies (CBC, ESR, autoimmune serologies) are obtained when systemic disease is suspected.
\end{itemize}

\section*{Treatment Options}
\begin{itemize}
  \item Topical therapies (corticosteroids, antifungals, antibiotics, emollients) are first-line for most localized skin conditions.
  \item Systemic therapy (oral antihistamines, antibiotics, antifungals, or immunosuppressants) is reserved for extensive or refractory cases.
  \item Phototherapy or biologic agents may be indicated for chronic inflammatory skin diseases (psoriasis, eczema).
  \item Referral to dermatology is appropriate for uncertain diagnoses, treatment failure, or skin lesions suspicious for malignancy.
\end{itemize}

\section*{References}
\begin{thebibliography}{99}
\bibitem{ref1} Bolognia JL, Schaffer JV, Cerroni L. \textit{Dermatology}, 4th ed. Elsevier, 2018.
\bibitem{ref2} Ferri FF. "Approach to the patient with rash." In: \textit{Ferri's Clinical Advisor 2024}. Elsevier, 2024:1205--1210.
\bibitem{ref3} Ely JW, Seabury Stone M. "The generalized rash: Part I. Differential diagnosis." \textit{American Family Physician}, 2010;81(6):726--734.
\bibitem{ref4} Wolff K, Johnson RA, Saavedra AP. \textit{Fitzpatrick's Color Atlas and Synopsis of Clinical Dermatology}, 8th ed. McGraw-Hill, 2017.
\bibitem{ref5} Drago LA, Davis MDP, Camilleri M, et al. "Clinical and histopathologic features of cutaneous drug eruptions." \textit{Journal of the American Academy of Dermatology}, 2020;83(4):1005--1019.
\bibitem{ref6} Lio PA, Bieber T. "Atopic dermatitis: Global epidemiology and risk factors." \textit{The Lancet}, 2020;396(10247):345--360.
\end{thebibliography}

\clearpage
